\begin{center}
{\Large\parbox{0.94\textwidth}{\centering\ThesisTitle}\par}
{\large\GroupName\par}
\vspace{2pt}
{\small\AuthorList\par}
{\large Degree \DegreeName\par}
{\large \InstitutionName\par}
{\large\ConvocationYear\par}
\end{center}

\vspace{0.25in}
\begin{center}
{\Large Abstract\par}
\end{center}
\addcontentsline{toc}{chapter}{Abstract}

This thesis presents an end-to-end fraud risk scoring system built on the
IEEE-CIS transaction dataset. Apache Spark performs typed ingestion, quality
auditing, chronological splitting, leakage-controlled preprocessing, feature
engineering, and Parquet export. A model pipeline compares logistic regression,
LightGBM, XGBoost, and CatBoost using ROC-AUC and precision--recall AUC, then
packages the selected model with categorical mappings and SHAP explanations.
FastAPI exposes scoring, transaction exploration, imports, dataset loading, and
human review, while a React dashboard supports investigation. The current
official data contract reports 94 passed verification checks. The official V2
CatBoost balanced package reports validation ROC-AUC 0.9081210 and PR-AUC
0.5805993, with holdout ROC-AUC 0.8780755 and PR-AUC 0.4266435. Its isotonic
calibrator and thresholds are packaged, but the promotion gate returns
\texttt{not\_promoted}; the configured serving version remains unchanged.
Limitations include local-only MLflow tracking, duplicated threshold ownership,
offline--online historical-feature mismatch, heuristic fallback risk,
single-node Python training, in-memory jobs, and missing production security
controls. The result is a coherent academic demonstration with explicit
architecture gaps, not a production-readiness claim.
